\documentclass[12pt]{article}
\usepackage[utf8]{inputenc}
\usepackage[russian]{babel}
\usepackage{amsmath}
\usepackage{geometry}

\geometry{left=2cm, right=2cm, top=2cm, bottom=2cm}

\title{Признаки равенства треугольников}
\author{losharics}
\date{\today}

\begin{document}

\maketitle

\section{Что такое равные треугольники?}

Два треугольника называются равными, если их можно совместить при наложении. Это значит, что у них должны быть равны все стороны и все углы.

Если мы возьмем треугольник $\triangle LOX$ и треугольник $\triangle L_1O_1X_1$, то из их равенства ($\triangle LOX = \triangle L_1O_1X_1$) следует:
\begin{enumerate}
    \item Равенство сторон: $LO = L_1O_1$, $OX = B_1C_1$ (то есть $O_1X_1$), $XL = X_1L_1$.
    \item Равенство углов: $\angle L = \angle L_1$, $\angle O = \angle O_1$, $\angle X = \angle X_1$.
\end{enumerate}

Чтобы доказать, что треугольники равны, не нужно проверять все 6 условий. Достаточно проверить только 3 главных элемента. Для этого есть три признака.

\section{Три главных признака равенства}

\subsection{Первый признак (по двум сторонам и углу между ними)}
Если две стороны и угол между ними одного треугольника равны двум сторонам и углу между ними другого треугольника, то такие треугольники равны.

Для наших треугольников это выглядит так:
\[
\begin{cases}
LO = L_1O_1 \\
LX = L_1X_1 \\
\angle L = \angle L_1
\end{cases}
\implies \triangle LOX = \triangle L_1O_1X_1
\]

\subsection{Второй признак (по стороне и двум углам)}
Если сторона и два прилежащих к ней угла одного треугольника равны стороне и двум углам другого треугольника, то они равны.

Формула:
\[
\begin{cases}
LO = L_1O_1 \\
\angle L = \angle L_1 \\
\angle O = \angle O_1
\end{cases}
\implies \triangle LOX = \triangle L_1O_1X_1
\]

\subsection{Третий признак (по трем сторонам)}
Если все три стороны одного треугольника равны трем сторонам другого треугольника, то эти треугольники равны.

Формула:
\[
\begin{cases}
LO = L_1O_1 \\
OX = O_1X_1 \\
XL = X_1L_1
\end{cases}
\implies \triangle LOX = \triangle L_1O_1X_1
\]

\section{Прямоугольные треугольники}
Для прямоугольных треугольников все еще проще, потому что один угол всегда равен $90^\circ$. Их можно доказать по двум элементам:
\begin{itemize}
    \item По двум катетам
    \item По катету и острому углу
    \item По гипотенузе и острому углу
    \item По гипотенузе и катету
\end{itemize}

\section*{! Важно знать}
Нельзя доказать равенство треугольников только по трем углам ($\angle L, \angle O, \angle X$). Если углы одинаковые, треугольники могут быть просто похожими (подобными), но иметь разный размер (например, один маленький, а другой огромный).

\end{document}